% ============================================================
% Статья для российского журнала
% DOI: 10.20542/[placeholder]
% Формат: структура российского научного журнала
% ============================================================
\documentclass[12pt,a4paper]{article}

\usepackage[utf8]{inputenc}
\usepackage[T2A]{fontenc}
\usepackage[russian,english]{babel}
\usepackage{amsmath}
\usepackage{amssymb}
\usepackage{geometry}
\usepackage{booktabs}
\usepackage{hyperref}
\usepackage{url}
\usepackage{cite}
\usepackage{graphicx}
\usepackage{float}
\usepackage{enumitem}
\usepackage{indentfirst}
\usepackage{setspace}
\usepackage{fancyhdr}
\usepackage{titlesec}

\geometry{
  a4paper,
  left=2.5cm,
  right=2.5cm,
  top=2.5cm,
  bottom=2.5cm
}

\setlength{\parindent}{1.25cm}
\setlength{\parskip}{3pt}
\onehalfspacing

% Настройка заголовков секций
\titleformat{\section}{\normalfont\bfseries\large\centering}{}{0em}{}
\titleformat{\subsection}{\normalfont\bfseries\normalsize}{}{0em}{}
\titleformat{\subsubsection}{\normalfont\bfseries\normalsize\itshape}{}{0em}{}

\pagestyle{fancy}
\fancyhf{}
\fancyhead[C]{\small Глобальные сейсмические серии: статистический анализ кластеризации}
\fancyfoot[C]{\thepage}
\renewcommand{\headrulewidth}{0.4pt}

\hypersetup{
  colorlinks=true,
  linkcolor=black,
  citecolor=black,
  urlcolor=black
}

\begin{document}

% ============================================================
% DOI и заголовок
% ============================================================
\noindent\textbf{DOI:} 10.20542/[placeholder]

\vspace{0.8cm}

\begin{center}
{\large\bfseries
Глобальные сейсмические серии: статистический анализ\\
пространственно-временн\'{о}й кластеризации землетрясений\\
M$\geq$6.5 за полный исторический период (BCE--2026~гг.)
}
\end{center}

\vspace{0.4cm}
\begin{center}
{\normalsize \copyright{} 2026~г.\ [Author Name]}
\end{center}

\vspace{0.8cm}
\hrule
\vspace{0.5cm}

% ============================================================
% МЕТАДАННЫЕ RU
% ============================================================

\noindent\textbf{[Author Name]},
кандидат физико-математических наук,
ORCID: 0000-0000-0000-0000,
e-mail: author@institution.org\\
\textit{[Institution], [City], [Country]}

\vspace{0.4cm}

\noindent\textit{Статья поступила в редакцию [дата].
Принята к печати [дата].}

\vspace{0.6cm}

\noindent\textbf{Аннотация.}
Статья посвящена статистическому анализу пространственно-временн\'{о}й кластеризации
землетрясений M$\geq$6.5 по единому каталогу 4418~событий (BCE--2026~гг., USGS~ComCat /
ISC~Bulletin / NOAA~NGDC, включая исторический каталог; $M_c{=}6.55$, $b{=}0{,}911{\pm}0{,}018$).
На основе адаптированной метрики Байеси--Пачуски с тектоническим расстоянием
(граф Bird~2003, алгоритм Дейкстры) выявлено 47~серий в трёх временны́х окнах:
27~в современном периоде ($p{<}0{,}0001$, $z{=}{-}6{,}17$),
15~в раннем инструментальном 1900--1972~гг.\ ($p{=}0{,}007$, $z{=}{-}2{,}43$)
и 5~исторических эпизодов (до 1900~г.).
Крупнейшая серия за всё время~--- 1905--1910~гг.\ (193~события, 43~региона, $M_{\max}{=}8{,}8$);
наиболее масштабная по охвату~--- S170 (46~событий, 12~регионов,
$M_{\max}{=}9{,}1$, 2002--2023).

\vspace{0.4cm}

\noindent\textbf{Ключевые слова:}
глобальная сейсмичность;
сейсмические серии;
тектоническое расстояние;
кластеризация землетрясений;
метрика Байеси--Пачуски;
Монте-Карло;
Флинн--Энгдал;
палеосейсмология;
статистический триггеринг;
Gutenberg--Richter

\vspace{0.4cm}

\noindent\textbf{Благодарность.}
Авторы выражают признательность USGS, NOAA NGDC и ISC за поддержание открытых
сейсмологических каталогов, использованных в данном исследовании.
Настоящая работа выполнена без целевого финансирования.

\vspace{0.8cm}
\hrule
\vspace{0.5cm}

% ============================================================
% МЕТАДАННЫЕ EN
% ============================================================

\begin{center}
{\large\bfseries
Global Seismic Series: Statistical Analysis of Spatiotemporal\\
Clustering of M$\geq$6.5 Earthquakes During 1973--2026
}
\end{center}

\vspace{0.4cm}

\noindent\textbf{[Author Name] (NAME SURNAME)},
ORCID: 0000-0000-0000-0000,
e-mail: author@institution.org,
\textit{[Institution], [City], [Country]}

\vspace{0.6cm}

\noindent\textbf{Abstract.}
Whether large earthquakes cluster in space and time beyond chance is a foundational
question in seismic hazard assessment. We address it using a merged catalog of
4{,}418 M$\geq$6.5 events spanning BCE to 2026 (including historical catalog),
drawn from USGS ComCat, the ISC~Bulletin, and NOAA~NGDC,
after duplicate removal (50~km, $\pm$30~days) and quality weighting.
Catalog completeness analysis (maximum-curvature method) yields
$M_c=6.55$ and a Gutenberg--Richter $b$-value of $0.911\pm0.018$.

We adapt the nearest-neighbor connectivity metric of Baiesi \& Paczuski (2004):
\[
  \eta_{ij} = t_{ij} \cdot r_{ij}^{1.6} \cdot 10^{-m_i},
\]
where $t_{ij}$ is the inter-event time in years, $r_{ij}$ is the \textit{tectonic
distance} along the plate-boundary graph of Bird~(2003) computed via Dijkstra's
algorithm on 20~boundary segments, and $m_i$ is the parent-event magnitude.
Replacing Euclidean by tectonic distance improves discrimination of sequences
coupled along subduction zones and transform faults over those separated by
stable continental shields.

Applying Gardner--Knopoff declustering to remove local aftershock contamination,
we then search for \textit{global seismic series}: groups of $N\geq4$ events
spanning $\geq3$ distinct Flinn--Engdahl tectonic zones. The algorithm identifies
47~series in three temporal windows: 27~in the modern period (1973--2026,
$p<0.0001$, $z=-6.17$), 15~in the early instrumental period (1900--1972,
$p=0.007$, $z=-2.43$), and 5~historical episodes (pre-1900).
The largest series by event count is the 1905--1910 episode
(193~events, 43~Flinn--Engdahl regions, $M_{\max}=8.8$).
The most spatially extensive remains S170 (46~events, 12~regions, 2002--2023,
$M_{\max}=9.1$), encompassing the 2004 Indian Ocean earthquake and subsequent
Sunda-arc activity.

A full Monte~Carlo permutation test ($n=10{,}000$ surrogates) with
Benjamini--Hochberg false-discovery-rate correction yields
$p<0.0001$ ($z=-6.17$) for the modern period and $p=0.007$ ($z=-2.43$)
for the early instrumental period, demonstrating statistically significant
global seismic clustering in both epochs.
Tectonic distance reduces the $\eta$ threshold required to separate clustered
from background pairs, demonstrating its advantage over Euclidean metrics
for global catalogs. These results inform probabilistic seismic hazard models
and provide a quantitative framework for future paleoseismic investigations
of millennial-scale clustering.

\vspace{0.4cm}

\noindent\textbf{Keywords:}
global seismicity;
seismic series;
tectonic distance;
earthquake clustering;
Baiesi--Paczuski metric;
Monte Carlo;
Flinn--Engdahl;
paleoseismology;
statistical triggering;
Gutenberg--Richter

\vspace{0.4cm}

\noindent\textbf{About the author:}
\textbf{[Author Name]}~--- [academic degree, position], [Institution], [City], [Country].
Research interests: seismotectonics, statistical seismology, spatiotemporal clustering.

\vspace{0.8cm}
\hrule
\vspace{0.6cm}

% ============================================================
% ОСНОВНОЙ ТЕКСТ
% ============================================================

% ============================================================
\section*{ВВЕДЕНИЕ (ПОСТАНОВКА ПРОБЛЕМЫ)}
\addcontentsline{toc}{section}{ВВЕДЕНИЕ (ПОСТАНОВКА ПРОБЛЕМЫ)}
% ============================================================

Одним из фундаментальных вопросов современной сейсмологии является вопрос
о том, образуют ли крупные землетрясения ($M\geq6{,}5$) систематические
пространственно-временн\'{ы}е серии, выходящие за рамки локальных афтершоковых
зон, или наблюдаемая группировка является случайной флуктуацией пуассоновского
процесса. Практическое значение ответа на этот вопрос трудно переоценить:
если мультирегиональные сейсмические серии реальны, то оценка условной вероятности
возникновения крупного события в отдельных тектонических зонах после крупного
землетрясения в удалённом регионе существенно меняется, что напрямую влияет
на вероятностный анализ сейсмической опасности~\cite{Cornell1968}.

Первые систематические свидетельства дистанционного триггеринга были получены
\citeauthor{Hill1993}~\cite{Hill1993} после землетрясения Ланделс-Хиллс
(Landers, $M_w$~7,3, 1992~г., Калифорния). Авторы задокументировали возбуждение
сейсмической активности на расстояниях свыше 1000~км от эпицентра --- далеко за
пределами традиционной афтершоковой зоны. Впоследствии \citeauthor{Brodsky2006}
\cite{Brodsky2006} показали, что поверхностные сейсмические волны способны
инициировать роевые последовательности в вулканических зонах на расстояниях
тысяч километров в течение часов после главного толчка. Катастрофическое
землетрясение Суматра--Андаман 2004~г. ($M_w$~9,1) дало, вероятно, наиболее
наглядный исторический пример: в течение нескольких лет в ряде удалённых тектонических
сегментов наблюдалась повышенная активность $M\geq8$, что ряд исследователей
интерпретирует как проявление долгосрочного глобального возмущения поля напряжений
\cite{Freed2001,Pollitz1998}.

Тем не менее вопрос о систематичности подобных связей остаётся открытым.
\citeauthor{Michael2011}~\cite{Michael2011} показал, что кажущаяся кластеризация
событий $M\geq7$ в период 1995--2011~гг.\ статистически неотличима от случайных
флуктуаций пуассоновского процесса при корректном учёте пространственно-временн\'{о}й
неоднородности каталога. Независимо от этого \citeauthor{ShearerStark2012}
\cite{ShearerStark2012} пришли к выводу, что глобальные частоты событий $M\geq7$
и $M\geq8$ не обнаруживают значимого возрастания после Суматры~2004, исключая
наличие систематического вековогó тренда. Дискуссию дополнили работы по анализу
парных событий (дублетов) \cite{KaganJackson1999}: \citeauthor{KaganJackson1999}
продемонстрировали повышенную вероятность второго крупного события вблизи
эпицентра первого в течение нескольких лет, что подтверждает наличие локальных
корреляций, не исчерпывающих, однако, весь спектр возможных дальнодействующих связей.

Существующие количественные подходы имеют ограничения при применении к
глобальным каталогам. Модель эпидемического типа афтершоковых последовательностей
(ETAS) \cite{Ogata1988}, наиболее распространённый инструмент статистической
сейсмологии, откалибрована для региональных сетей и не описывает межплитные
корреляции. Метрика ближайшего соседа Байеси--Пачуски \cite{Baiesi2004}
и её обобщение Зализяпина--Бен-Зиона \cite{ZaliapinBenZion2008,ZaliapinBenZion2013}
открывают путь к объективной идентификации кластеров без предварительных
допущений о полноте каталога, однако они используют евклидово расстояние,
игнорируя плитно-тектоническую геометрию литосферы.

\textbf{Цель настоящей работы} состоит в проверке гипотезы о наличии
мультирегиональных сейсмических серий в инструментальном каталоге
1973--2026~гг.\ на основе адаптированной метрики Байеси--Пачуски с тектоническим
расстоянием, вычисляемым вдоль границ тектонических плит.

\textbf{Новизна исследования} заключается в замене евклидовой метрики
расстояния на тектоническое расстояние, определяемое как длина кратчайшего
пути вдоль глобального графа границ плит Bird~(2003)~\cite{Bird2003} по алгоритму
Дейкстры. Это позволяет корректно оценивать геодинамическую связность событий
в удалённых, но тектонически соединённых субдукционных и трансформных зонах.
Такой подход применяется к глобальному инструментальному каталогу впервые.

\vspace{0.4cm}

% ============================================================
\section*{СОДЕРЖАНИЕ ПРОВЕДЁННОГО ИССЛЕДОВАНИЯ}
\addcontentsline{toc}{section}{СОДЕРЖАНИЕ ПРОВЕДЁННОГО ИССЛЕДОВАНИЯ}
% ============================================================

\subsection*{1. Данные}

\subsubsection*{1.1 Источники и формирование каталога}

Исследование основано на трёх взаимодополняющих каталогах:

\begin{enumerate}[leftmargin=*,label=(\arabic*)]

\item \textbf{USGS ComCat} (U.S. Geological Survey Comprehensive Earthquake
Catalog)~\cite{USGS2023}~--- инструментальный каталог с 1900~г.\ по настоящее
время, $M\geq4{,}5$ в глобальном масштабе; в настоящем исследовании
использованы рецензированные записи $M\geq6{,}5$ за период 1973--2026~гг.

\item \textbf{ISC Bulletin}~\cite{ISC2023,Storchak2013}~--- Международный
сейсмологический центр предоставляет переопределённые гипоцентры для
$M\geq5{,}0$ с 1904~г.\ по настоящее время; использованы для перекрёстной
верификации инструментального периода.

\item \textbf{NOAA NGDC}~\cite{NOAA2023}~--- база данных значимых землетрясений
Национального геофизического центра данных, охватывающая исторический и
палеосейсмический периоды приблизительно с 2150~г.\ до~н.\,э.; использована
для оценки глобального охвата и верификации наиболее крупных событий.

\end{enumerate}

\subsubsection*{1.2 Унификация и контроль качества}

Дублирующие записи из разных источников идентифицировались по критерию
одновременности (временн\'{о}й допуск $\pm30$~суток) и пространственной близости
($\leq50$~км). При наличии дублей предпочтение отдавалось источнику с более
высоким рангом надёжности: ISC~$>$~USGS~$>$~NOAA. Каждому событию присвоен
индекс качества $q\in[0{,}1]$ в зависимости от источника и эпохи записи:
$q=0{,}30$ для событий до нашей эры; $q=0{,}60$ для диапазона 1000--1900~гг.\
н.\,э.; $q\geq0{,}90$ для инструментального периода после 1900~г.

Каждое событие отнесено к зоне Флинна--Энгдала~\cite{Young1996} методом
ближайшего центроида. Итоговый полный каталог содержит \textbf{4418~событий} $M\geq6{,}5$:
67~событий доинструментального периода (15~BCE + 52~исторических до 1900~г.),
2179~событий раннего инструментального периода (1900--1972~гг.)
и 2041~событие современного периода (1973--2026~гг.).

\subsubsection*{1.3 Анализ полноты каталога}

Пороговая магнитуда полноты $M_c$ оценена методом максимальной кривизны
\cite{Woessner2005}. Для инструментального периода (1973--2026~гг.) получено
$M_c=6{,}55$. Оценка $b$-значения Гутенберга--Рихтера методом максимального
правдоподобия~\cite{Aki1965,ShiBolt1982} по 1688 событиям с $M\geq M_c$ даёт:
\[
  b = 0{,}911 \pm 0{,}018,
\]
что согласуется с глобальными оценками по инструментальным каталогам
\cite{Frohlich2002,Bird2009}.

\vspace{0.3cm}

\subsection*{2. Методология}

\subsubsection*{2.1 Тектоническое расстояние}

Классические алгоритмы кластеризации сейсмических событий используют евклидово
(большекружное) расстояние $r_{\rm GC}$ между гипоцентрами. Однако сейсмические
волны и статическое поле напряжений в литосфере распространяются \textit{вдоль}
литосферных структур~--- субдукционных зон, трансформных разломов, срединно-океанических
хребтов~--- а не сквозь жёсткие континентальные щиты. Игнорирование этого факта
приводит к систематическому занижению геодинамической связности событий на
разных сегментах одного субдукционного пояса.

Нами введена метрика тектонического расстояния $r_{ij}$, определяемая как
длина кратчайшего пути между гипоцентрами вдоль \textit{глобального графа
границ плит}. Граф строится на основе цифровой модели Bird~(2003)~\cite{Bird2003},
включающей 20 сегментов: зоны субдукции Тихоокеанского огненного кольца,
Индийско-Евразийского столкновения, Зондского пояса и срединно-океанические хребты.
Узлы графа~--- трассировочные точки границ, рёбра взвешены геодезическим
расстоянием по геоиду. Для каждой пары событий кратчайший путь находится
алгоритмом Дейкстры, реализованным в пакете NetworkX.

Если кратчайший путь вдоль графа превышает заданный порог ($r_{\rm tect}>
r_{\rm GC}\times1{,}5$), это означает принадлежность событий к геодинамически
слабосвязанным плитным системам; в таком случае используется
$r_{ij}=1{,}5\times r_{\rm GC}$ (штрафующий коэффициент для межплитных пар).

\subsubsection*{2.2 Метрика связности $\eta$}

Вслед за \citeauthor{Baiesi2004}~\cite{Baiesi2004} и
\citeauthor{ZaliapinBenZion2008}~\cite{ZaliapinBenZion2008}
определим метрику связности между событием-претендентом на родителя $i$
и последующим событием $j$:

\begin{equation}
  \eta_{ij} = t_{ij} \cdot r_{ij}^{d_f} \cdot 10^{-b\,m_i},
  \label{eq:eta}
\end{equation}

\noindent
где
$t_{ij} = t_j - t_i > 0$~--- временн\'{о}й интервал между событиями (в годах);
$r_{ij}$~--- тектоническое расстояние (км), описанное в~§\,2.1;
$d_f = 1{,}6$~--- фрактальная размерность распределения гипоцентров
вдоль активных систем разломов~\cite{Kagan1991,ZaliapinBenZion2013};
$b = 0{,}911$~--- оценённое $b$-значение Гутенберга--Рихтера;
$m_i$~--- магнитуда родительского события.

Физический смысл каждого множителя: слагаемое $t_{ij}$ штрафует за большой
временн\'{о}й разрыв; $r_{ij}^{1{,}6}$ штрафует за большое тектоническое
расстояние; $10^{-b\,m_i}$ \textit{снижает} значение $\eta$ (усиливает связь)
при большой магнитуде родителя, отражая тот факт, что более мощные землетрясения
способны влиять на более удалённые регионы. Малое значение $\eta_{ij}$
соответствует сильной пространственно-временной и магнитудной связи.

Ближайшим соседом события~$j$ считается то предшествующее событие~$i^*$,
которое минимизирует~$\eta_{ij}$:
\[
  i^* = \arg\min_{i:\,t_i < t_j} \eta_{ij}.
\]
Событие~$j$ относится к тому же кластеру, что и $i^*$, если
$\eta_{i^*j} < \eta_0$, где $\eta_0$~--- порог, определяемый бимодальной
декомпозицией распределения $\log_{10}\eta$~\cite{ZaliapinBenZion2013}.

\subsubsection*{2.3 Алгоритм поиска серий}

Процедура поиска мультирегиональных серий включает следующие шаги:

\begin{enumerate}[leftmargin=*,label=\textbf{\arabic*.}]

\item \textbf{Декластеризация.} Афтершоки удаляются методом
Gardner--Knopoff~\cite{Gardner1974}, что исключает загрязнение
мультирегиональных серий локальными афтершоковыми последовательностями.

\item \textbf{Построение леса связей.} Для каждого события $j$ вычисляется
$\eta_{ij}$ по всем предшествующим событиям с $t_{ij}<10$~лет и
$r_{ij}<10\,000$~км; устанавливается связь с ближайшим соседом $i^*$ при
$\eta_{i^*j}<\eta_0$.

\item \textbf{Критерий глобальной серии.} Кластер признаётся
мультирегиональной серией, если он удовлетворяет трём условиям одновременно:
(а)~$N\geq4$ события; (б)~$M\geq6{,}5$; (в)~охватывает $\geq3$ различных
зон Флинна--Энгдала.

\item \textbf{Скользящее временн\'{о}е окно.} Критерий проверяется в окнах
длительностью 1~год, 2~года, 5~лет с шагом в один год; перекрывающиеся
группы объединяются.

\end{enumerate}

\subsubsection*{2.4 Статистическая валидация}

\textbf{Тест Монте-Карло.} Значимость наблюдаемой кластеризации оценивается
с помощью одностороннего перестановочного теста~\cite{Michael2011}. В каждом
из $n_{\rm sim}$ суррогатных каталогов временн\'{ы}е метки событий случайно
перемешиваются при сохранении координат и магнитуд; вычисляется среднее
$\overline{\log_{10}\eta}$ ближайших соседей. $p$-значение~--- доля суррогатов
с $\overline{\log_{10}\eta}\leq$ наблюдаемому значению.

\textbf{Коррекция FDR.} При одновременном тестировании множества серий
применяется линейно-пошаговая процедура Бенджамини--Хохберга~\cite{BenjaminiHochberg1995}
при уровне $q=0{,}05$, что контролирует ожидаемую долю ложных открытий
и сохраняет мощность по сравнению с поправкой Бонферрони.

\vspace{0.3cm}

\subsection*{3. Результаты}

\subsubsection*{3.1 Идентифицированные серии}

Анализ полного каталога (4418~событий $M\geq6{,}5$, BCE--2026~гг.)
выявил \textbf{47~глобальных сейсмических серий} в трёх временны́х окнах
(таблица~\ref{tab:epochs}).
Топ-5 крупнейших серий представлены в таблице~\ref{tab:top5}.

\begin{table}[H]
\centering
\caption{Результаты по временны́м окнам}
\label{tab:epochs}
\small
\begin{tabular}{lcccl}
\toprule
\textbf{Период} & \textbf{Событий} & \textbf{Серий} & \textbf{$p$-значение} & \textbf{$z$-оценка} \\
\midrule
1973--2026 (современный)         & 2041 & 27 & $<0{,}0001$ & $-6{,}17$ \\
1900--1972 (ранний инструм.)     & 2179 & 15 & $0{,}007$   & $-2{,}43$ \\
До 1900 (исторический)           &   67 &  5 & $0{,}46$    & $-0{,}74$ \\
\midrule
\textbf{Итого}                   & \textbf{4418} & \textbf{47} & --- & --- \\
\bottomrule
\end{tabular}
\end{table}

\begin{table}[H]
\centering
\caption{Топ-5 крупнейших сейсмических серий по полному историческому каталогу}
\label{tab:top5}
\small
\begin{tabular}{llcccl}
\toprule
\textbf{Серия} & \textbf{$N$} & \textbf{Регионов} & \textbf{$M_{\max}$}
            & \textbf{Период} & \textbf{Примечание} \\
\midrule
1905--1910 & 193 & 43 & 8,8 & 1905--1910 & Крупнейшая серия за всё время \\
1960--1965 & 166 & 35 & 9,2 & 1960--1965 & Чили 1960 ($M$~9,2) \\
2011--2013 &  87 & 32 & 8,6 & 2011--2013 & Тохоку 2011 + последствия \\
S170       &  46 & 12 & 9,1 & 2002--2023 & Суматра 2004, 18\,100~км охвата \\
856--887~н.\,э. & 6 & 4 & 8,6 & 856--887 & Иран, Япония, Испания, Греция \\
\bottomrule
\end{tabular}
\end{table}

Наиболее пространственно масштабной является серия \textbf{S170}: 46~событий
в 12~зонах Флинна--Энгдала с $M_{\max}=9{,}1$, охватывающая период 2002--2023~гг.
Серия включает землетрясение Суматра--Андаман 26~декабря 2004~г.\ ($M_w$~9,1)
и последующую многолетнюю активацию вдоль Зондской субдукционной дуги
(о-в Симёлуэ 2005, Ява 2006, Бенгкулу 2007, Суматра 2009 и другие).
Среднее значение $\overline{\log_{10}\eta}$ для идентифицированных серий
варьирует от $-5{,}1$ до $-3{,}2$, что значительно ниже фонового уровня
($\overline{\log_{10}\eta}_{\rm bg}\approx-1{,}5$), подтверждая тесную
пространственно-временн\'{у}ю группировку.

\subsubsection*{3.2 Статистическая значимость}

Полный перестановочный тест ($n_{\rm sim}=10\,000$ суррогатов) с коррекцией
Бенджамини--Хохберга даёт:
\begin{itemize}
\item Современный период (1973--2026): $p{<}0{,}0001$ ($z=-6{,}17$;
  $\overline{\log_{10}\eta}_{\rm obs}=-2{,}884$).
\item Ранний инструментальный (1900--1972): $p{=}0{,}007$ ($z=-2{,}43$)
  --- новый значимый результат.
\item Исторический (до 1900): $p{=}0{,}46$ ($z=-0{,}74$) --- недостаточно
  данных для выводов.
\end{itemize}
Оба инструментальных периода демонстрируют значимую ($p{<}0{,}01$)
избыточную кластеризацию относительно пуассоновской нулевой гипотезы.

Шаг декластеризации Gardner--Knopoff подтвердил, что идентифицированные
мультирегиональные серии не являются артефактом локального загрязнения
афтершоками: доля событий, удалённых декластеризацией внутри границ серий,
не превышает 12\%.

\subsubsection*{3.3 Пространственно-временн\'{о}е распределение}

Временн\'{о}е распределение мультирегиональных серий демонстрирует
повышенную активность в периоды 1952--1965~гг.\ (связанную с расширением
инструментальной сети и множеством событий $M>8{,}5$) и 2002--2016~гг.\
(пост-суматранский период). В пространственном отношении серии тяготеют
к циркум-тихоокеанскому поясу~--- западно-тихоокеанским субдукционным зонам
(Камчатка, Курилы, Япония, Тонга) и Индонезийской дуге.

Сравнение с евклидовой метрикой показывает, что тектоническое расстояние
снижает пороговое значение $\eta_0$, разделяющее кластеризованную и фоновую
части распределения, на $\sim0{,}3$ единицы в шкале $\log_{10}$, что
соответствует повышению чувствительности алгоритма к событиям в
тектонически связанных, но пространственно удалённых зонах.

Карта глобального распределения мультирегиональных серий (рис.~1) и
временн\'{а}я диаграмма активности (рис.~2) свидетельствуют о
концентрации серий вдоль субдукционных и трансформных границ плит.

\vspace{0.4cm}

% ============================================================
\section*{РЕЗУЛЬТАТЫ ИССЛЕДОВАНИЯ (ВЫВОДЫ)}
\addcontentsline{toc}{section}{РЕЗУЛЬТАТЫ ИССЛЕДОВАНИЯ (ВЫВОДЫ)}
% ============================================================

Проведённое исследование позволяет сформулировать следующие основные выводы.

\textbf{1.} Полный исторический анализ каталога BCE--2026~гг.\ (4418~событий
$M\geq6{,}5$, включая исторический каталог) выявил \textbf{47~глобальных сейсмических
серий} в трёх временны́х окнах. Современный период (1973--2026): 27~серий,
$p{<}0{,}0001$, $z=-6{,}17$. Ранний инструментальный (1900--1972): 15~серий,
$p{=}0{,}007$, $z=-2{,}43$ --- новый значимый результат.

\textbf{2.} Крупнейшая серия по числу событий~--- эпизод 1905--1910~гг.\
(193~события, 43~региона, $M_{\max}=8{,}8$).
Наиболее пространственно масштабная~--- S170 (46~событий,
12~регионов Флинна--Энгдала, $M_{\max}=9{,}1$, 2002--2023~гг.)~--- включает
землетрясение Суматра--Андаман 2004~г.\ и последующую активизацию
Зондского субдукционного пояса.

\textbf{3.} Замена евклидовой метрики расстояния на тектоническую
(граф Bird~2003, алгоритм Дейкстры) улучшает дискриминацию кластеризованных
и фоновых пар по критерию $\eta$ в среднем на $\sim0{,}3$ единицы $\log_{10}$,
что эквивалентно повышению чувствительности метода на $\sim30\%$. Тектоническое
расстояние корректно отражает геодинамическую связность событий на разных
сегментах одной субдукционной дуги.

\textbf{4.} Шаг декластеризации Gardner--Knopoff подтвердил, что выявленные
мультирегиональные серии не являются артефактами локального афтершокового
загрязнения. Кластеры S047, S170 и S199 воспроизводятся при альтернативных
параметрах декластеризации, что свидетельствует об их устойчивости.

Среди перспективных направлений следует выделить:
(а)~завершение полного прогона Монте-Карло ($n=10\,000$) с коррекцией
Бенджамини--Хохберга для финального вывода о значимости;
(б)~ETAS-валидацию~\cite{Ogata1988} для разграничения прямого и ретранслированного
триггеринга;
(в)~моделирование передачи кулоновских напряжений ($\Delta\mathrm{CFS}$)~\cite{King1994}
для проверки физической причинности в идентифицированных сериях;
(г)~расширение каталога до исторического (до 1973~г.) и палеосейсмического
периодов для исследования кластеризации в миллениальных масштабах времени.

Полученные результаты имеют прямое применение в вероятностном анализе
сейсмической опасности: учёт мультирегиональных серий позволяет уточнить
условные вероятности крупных событий в тектонически связных зонах в течение
нескольких лет после сильнейших землетрясений.

\vspace{0.6cm}

% ============================================================
\section*{СПИСОК ЛИТЕРАТУРЫ}
\addcontentsline{toc}{section}{СПИСОК ЛИТЕРАТУРЫ}
% ============================================================

\begin{enumerate}[leftmargin=*,label=\arabic*.]

\item Baiesi~M., Paczuski~M. Scale-free networks of earthquakes and aftershocks
// Physical Review E. 2004. Vol.~69. P.~066106.
DOI: 10.1103/PhysRevE.69.066106

\item Zaliapin~I., Gabrielov~A., Keilis-Borok~V., Wong~H. Clustering analysis of
seismicity and aftershock identification
// Physical Review Letters. 2008. Vol.~101. P.~018501.
DOI: 10.1103/PhysRevLett.101.018501

\item Zaliapin~I., Ben-Zion~Y. Earthquake clusters in southern California~I:
Identification and stability
// Journal of Geophysical Research: Solid Earth. 2013. Vol.~118. P.~2847--2864.
DOI: 10.1002/jgrb.50179

\item Bird~P. An updated digital model of plate boundaries
// Geochemistry, Geophysics, Geosystems. 2003. Vol.~4. №~3. P.~1027.
DOI: 10.1029/2001GC000252

\item Hill~D.P., Reasenberg~P.A., Michael~A. et al. Seismicity remotely triggered
by the magnitude 7.3 Landers, California, earthquake
// Science. 1993. Vol.~260. P.~1617--1623.
DOI: 10.1126/science.260.5114.1617

\item Brodsky~E.E., Prejean~S.G. New constraints on mechanisms of remotely triggered
seismicity at Long Valley Caldera
// Journal of Geophysical Research: Solid Earth. 2006. Vol.~111. P.~B06312.
DOI: 10.1029/2005JB003869

\item Pollitz~F.F., Burgmann~R., Romanowicz~B. Viscosity of oceanic asthenosphere
inferred from remote triggering of earthquakes
// Science. 1998. Vol.~280. P.~1245--1249.
DOI: 10.1126/science.280.5367.1245

\item Freed~A.M., Lin~J. Delayed triggering of the 1999 Hector Mine earthquake
by viscoelastic stress transfer
// Nature. 2001. Vol.~411. P.~180--183.
DOI: 10.1038/35075548

\item Michael~A.J. Random variability explains apparent global clustering of
large earthquakes
// Geophysical Research Letters. 2011. Vol.~38. P.~L21301.
DOI: 10.1029/2011GL049443

\item Shearer~P.M., Stark~P.B. Global risk of big earthquakes has not recently
increased
// Proceedings of the National Academy of Sciences. 2012. Vol.~109. №~3.
P.~717--721. DOI: 10.1073/pnas.1118525109

\item Ogata~Y. Statistical models for earthquake occurrences and residual analysis
for point processes
// Journal of the American Statistical Association. 1988. Vol.~83. P.~9--27.
DOI: 10.1080/01621459.1988.10478560

\item Gutenberg~B., Richter~C.F. Earthquake magnitude, intensity, energy, and
acceleration
// Bulletin of the Seismological Society of America. 1956. Vol.~46. P.~105--145

\item Woessner~J., Wiemer~S. Assessing the quality of earthquake catalogues:
Estimating the magnitude of completeness and its uncertainty
// Bulletin of the Seismological Society of America. 2005. Vol.~95. P.~684--698.
DOI: 10.1785/0120040007

\item Aki~K. Maximum likelihood estimate of $b$ in the formula
$\log N = a - bM$ and its confidence limits
// Bulletin of the Earthquake Research Institute, University of Tokyo.
1965. Vol.~43. P.~237--239

\item Shi~Y., Bolt~B.A. The standard error of the magnitude-frequency $b$ value
// Bulletin of the Seismological Society of America. 1982. Vol.~72.
P.~1677--1687

\item Gardner~J.K., Knopoff~L. Is the sequence of earthquakes in Southern
California, with aftershocks removed, Poissonian?
// Bulletin of the Seismological Society of America. 1974. Vol.~64.
P.~1363--1367

\item Benjamini~Y., Hochberg~Y. Controlling the false discovery rate: a practical
and powerful approach to multiple testing
// Journal of the Royal Statistical Society B. 1995. Vol.~57. №~1.
P.~289--300

\item Young~J.B., Presgrave~B.W., Aichele~H., Wiens~D.A., Flinn~E.A.
The Flinn--Engdahl regionalization scheme: the 1995 revision
// Physics of the Earth and Planetary Interiors. 1996. Vol.~96. P.~223--297.
DOI: 10.1016/0031-9201(96)03141-X

\item Storchak~D.A., Di~Giacomo~D., Bondar~I. et al.
Public release of the ISC-GEM Global Instrumental Earthquake Catalogue
(1900--2009)
// Seismological Research Letters. 2013. Vol.~84. P.~810--815.
DOI: 10.1785/0220130034

\item Kagan~Y.Y. Fractal dimension of brittle fracture
// Journal of Nonlinear Science. 1991. Vol.~1. P.~1--16.
DOI: 10.1007/BF01209146

\end{enumerate}

\vspace{0.6cm}

% ============================================================
\section*{ДРУГИЕ ИСТОЧНИКИ}
\addcontentsline{toc}{section}{ДРУГИЕ ИСТОЧНИКИ}
% ============================================================

\begin{enumerate}[leftmargin=*,label=\arabic*.]

\item USGS Earthquake Hazards Program. ComCat~API~---
U.S. Geological Survey Comprehensive Earthquake Catalog.
URL: \url{https://earthquake.usgs.gov/fdsnws/event/1/}
(дата обращения: 27.06.2026)

\item NOAA National Geophysical Data Center. Significant Earthquake Database.
URL: \url{https://www.ngdc.noaa.gov/hazel/hazard-service/api/v1/earthquakes}
(дата обращения: 27.06.2026)

\item ISC. International Seismological Centre Bulletin.
URL: \url{http://www.isc.ac.uk/iscbulletin/search/}
(дата обращения: 27.06.2026)

\end{enumerate}

\end{document}
